\section{先完成一次小小的写作循环}
欢迎来到 Prims。你不必先学完所有 LaTeX 命令：先改一句话，编译一次，再试着让 AI 帮你推敲表达，就能开始使用。

\subsection{认识这个项目}
左侧文件树保存项目文件，中间是源代码，右侧是编译后的 PDF。小屏幕上可以通过“展开侧栏”和“显示预览”切换面板。打开主文件 \texttt{main.tex}，点击“编译”，这份教程就会出现在预览区。打开其他章节时，编译的仍是项目设置中的主文件。

\begin{center}
\begin{tabular}{ll}
\toprule
文件 & 用途 \\
\midrule
\texttt{main.tex} & 主文件：文档格式与章节入口 \\
\texttt{chapters/01-getting-started.tex} & 入门与第一轮练习 \\
\texttt{chapters/02-ai-writing.tex} & API 配置与选区 AI 对话 \\
\texttt{chapters/03-latex-examples.tex} & 公式、图片、表格与引用示例 \\
\texttt{references.bib} & 可替换的教学参考文献条目 \\
\texttt{figures/workflow.png} & 可编译、可供视觉模型读取的示例图片 \\
\texttt{README.md} & 不编译进 PDF 的速查说明 \\
\bottomrule
\end{tabular}
\end{center}

\subsection{试着修改这段话}
\begin{quote}
清晰的写作始于一个具体的问题。先写出你想表达的观点，再检查证据是否支持它，最后让句子变得简洁。
\end{quote}
\noindent 练习：把这段话中的“具体的问题”换成你的研究主题，点击“编译”，看看 PDF 的变化。编辑后的文档会自动保存；若网络断开或有协作冲突，请留意编辑器的状态提示。

\subsection{一些实用的位置}
“设置 → 个人资料”可以修改显示名称和界面语言；“设置 → 个人 AI 配置”可以填写或修改模型服务。管理员还会在设置中看到“站点管理”，用于管理邀请码和账号。你可以重命名这个欢迎项目，把它继续当作写作草稿使用。

\noindent 每个账号都有自己的欢迎项目。重复登录或刷新不会重复创建，删除后也不会自动重新出现。
